% ============================================================================
% Section 9: Comparison with Existing Theories
% ============================================================================

\section{Comparison with existing theories}
\label{sec:comparison}

UHM is not the only mathematical theory of consciousness. To situate it
precisely, this section provides a systematic comparison with the six
leading frameworks: Integrated Information
Theory (IIT 4.0)~\cite{iit4}, Global Neuronal Workspace Theory
(GNWT)~\cite{dehaene2011}, the Free Energy Principle (FEP)~\cite{friston2019},
Higher-Order Theories (HOT)~\cite{lau2011}, Hoffman and Prakash's Conscious Realism
and Conscious Agent Theory (CAT)~\cite{hoffman2014}, and Worden's recently proposed
Projective Wave Theory (PWT)~\cite{worden2026}.

\begin{figure}[H]
\centering
\input{figures/comparison-radar}
\caption{\textbf{Systematic comparison of UHM with leading theories of consciousness.}
Green = fully supported, yellow = partially addressed, red = absent. UHM is the only
theory that simultaneously provides a mathematical framework for the hard problem, derives physics, predicts critical
exponents, proves zombie impossibility, and provides an ethical framework.}
\label{fig:comparison}
\end{figure}

\subsection{Systematic comparison}

Table~\ref{tab:comparison} compares the theories across seven properties that a
complete theory of consciousness should possess. The expanded nine-row visual
comparison (hard problem, numerical threshold, critical exponents, falsifiable
predictions, physics derivation, zombie impossibility, consciousness-as-survival
criterion, computational tractability, and ethical framework) is reproduced in
Figure~\ref{fig:comparison}.

\begin{table}[H]
\centering
\caption{Systematic comparison of theories of consciousness. \checkmark = fully addresses;
$\sim$ = partially addresses; --- = does not address.}
\label{tab:comparison}
\small
\renewcommand{\arraystretch}{1.3}
\begin{tabular}{@{}p{2.9cm}ccccccc@{}}
\toprule
\textbf{Property} & \textbf{UHM} & \textbf{IIT 4.0} & \textbf{GNWT} & \textbf{FEP}
  & \textbf{HOT} & \textbf{CAT} & \textbf{PWT} \\
\midrule
Hard problem addressed &
  \checkmark & $\sim$ & --- & --- & --- & $\sim$ & $\sim$ \\
Numerical threshold &
  \checkmark\,(2/7) & ---\,($\Phi > 0$) & --- & --- & --- & --- & --- \\
Critical exponents &
  \checkmark\,($\beta{=}\tfrac{1}{4}$) & --- & --- & --- & --- & --- & --- \\
Falsifiable numerical predictions &
  22 & 0 & 0 & 0 & 0 & 0 & 0 \\
Computational complexity &
  $O(N^2)$ for $P$ & super-exp.\ for $\Phi$ & N/A & $O(N^2)$ for $F$
  & N/A & N/A & N/A \\
Derives physics &
  \checkmark\,(GR, QM, SM) & --- & --- & --- & --- & --- & --- \\
Zombie impossibility &
  \checkmark\,(theorem) & --- & --- & --- & --- & --- & --- \\
\bottomrule
\end{tabular}
\end{table}

\subsection{Theory-by-theory analysis}

\paragraph{IIT 4.0~\cite{iit4}.}
IIT is the closest competitor in formal ambition. It identifies five axioms of phenomenal
existence and derives postulates for physical substrates. Its central measure $\Phi$
(integrated information) quantifies how irreducible a system's cause--effect structure is.

However, IIT faces three problems that UHM resolves. First, \emph{computational
intractability}: computing $\Phi$ requires searching for the minimum information
partition over all partitions of the system, whose count grows super-exponentially
in the number of elements; practical full-$\Phi$ calculations are limited to roughly
ten to fifteen elements on current hardware, beyond which only approximation schemes
are available. UHM's purity measure $P = \Tr(\Gam^2)$ is $O(N^2) = O(49)$ by
comparison. Second, \emph{no numerical threshold}: IIT states that consciousness
corresponds to $\Phi > 0$, but any system with more than one element can have
$\Phi > 0$, including a simple grid of logic gates. UHM derives $\Pcrit = 2/7$ from
five independent arguments. Third, \emph{no explanation of why}: IIT postulates
that integrated information \emph{is} consciousness but does not explain why
irreducibility should feel like anything. UHM reframes this question via
two-aspect monism (Def.~\ref{def:two-aspect}): experience is not \emph{produced
by} $\Gam$; it is the internal aspect of $\Gam$.

The COGITATE adversarial collaboration (launched 2018, results published April 2025
in \emph{Nature}~\cite{cogitate2025}), funded by the Templeton World Charity Foundation's
\emph{Accelerating Research on Consciousness} initiative, tested IIT directly against
GNWT with $N = 256$ human participants across fMRI, MEG, and iEEG modalities. The
outcome was effectively a draw: IIT's prediction of sustained posterior synchronization
during conscious perception was not observed, while GNWT was challenged by the absence
of prefrontal ``ignition'' at stimulus offset.

\paragraph{GNWT~\cite{dehaene2011}.}
GNWT proposes that consciousness arises when information is ``ignited'' into a global
workspace of interconnected cortical neurons, enabling widespread access. The theory
successfully predicts the ``ignition'' pattern in EEG/fMRI --- a rapid, all-or-nothing
transition in neural activity when a stimulus crosses the conscious access threshold.

UHM shares GNWT's emphasis on avalanche dynamics (Prediction~16:
$T_{\mathrm{ign}} \sim (P - \Pcrit)^{-1}$), but adds two things GNWT lacks: (a)~a
specific numerical threshold ($\Pcrit = 2/7$), and (b)~an explanation of \emph{why}
broadcasting information should produce experience. In GNWT, the workspace is a functional
description; in UHM, the workspace is the E-sector of $\Gam$, and broadcasting is the
regeneration channel $\Regen$ whose strength is proportional to $\CohE$.

\paragraph{FEP~\cite{friston2019}.}
The Free Energy Principle provides a principled variational framework for self-organizing
systems. It unifies perception, action, and learning under the minimization of variational
free energy $F = D_{\mathrm{KL}}[q(\theta) \| p(\theta | o)] - \ln p(o)$. FEP is arguably
the most mathematically mature alternative to UHM.

The critical difference is that FEP does not distinguish conscious from unconscious free
energy minimization. A thermostat minimizes free energy. A rock at thermal equilibrium
minimizes free energy. FEP provides no threshold at which free energy minimization
transitions from mechanical to conscious. UHM resolves this by identifying the specific
structure ($\Gam \in \Dens(\mathbb{C}^7)$, the E-sector, $\CohE$) that makes
regeneration \emph{conscious} regeneration. Free energy minimization is necessary but
not sufficient; the Fano structure of self-observation channels is what distinguishes
a conscious system from a thermostat.

\paragraph{HOT~\cite{lau2011}.}
Higher-Order Theories propose that a mental state becomes conscious when it is the object
of a higher-order representation. UHM agrees that self-representation is necessary
(the self-modeling operator $\vp$, Section~\ref{subsec:phi-operator}), but adds:
(a)~a quantitative threshold $\Rth = 1/3$ for when self-modeling constitutes reflective
awareness; (b)~a proof that higher-order representations are bounded at depth 3
($\mathrm{SAD}_{\max} = 3$, Prediction~12); and (c)~an explanation of why
meta-representation feels like anything (two-aspect monism, not information processing).

\paragraph{Conscious Realism / Conscious Agent Theory (CAT)~\cite{hoffman2014}.}
Hoffman and Prakash's theory is the closest to UHM in ontological ambition:
both reject physicalism and place consciousness at the foundation. The 2014
paper defines a conscious agent as the 6-tuple
\begin{equation*}
  C \;=\; \bigl((X, \mathcal{X}),\; (G, \mathcal{G}),\; P,\; D,\; A,\; N\bigr)
\end{equation*}
where $(X, \mathcal{X})$ is the measurable space of experiences,
$(G, \mathcal{G})$ the measurable space of actions, and $P$, $D$, $A$ are
Markovian kernels (perception, decision, action) connecting these spaces to a
world $W$; $N$ counts message passes~\cite{hoffman2014}. The paper proves two
join theorems --- Undirected Join (Theorem~1) and Directed Join (Theorem~2)
--- which establish that conscious agents can be combined into new conscious
agents, the formal backbone of Hoffman's conscious-agent network construction.

However, CAT diverges from UHM in four critical respects:
\begin{enumerate}[leftmargin=2em]
  \item \textbf{No threshold for consciousness.} Every conscious agent is
        conscious by definition. UHM distinguishes L0 (formal interiority)
        through L2 (cognitive consciousness) via four computable thresholds.
        CAT cannot distinguish a thermostat from a human --- both are
        ``conscious agents'' by Definition~1.
  \item \textbf{Spacetime is illusion vs.\ emergent.} Within Hoffman's
        programme, the ``Fitness Beats Truth'' result~\cite{prakash2021fbt}
        argues that natural selection favours adaptive perception over
        veridical mapping of reality, motivating the claim that spacetime is
        a species-specific interface. UHM instead derives spacetime as
        \emph{emergent and physical}: $M^4 = \mathbb{R} \times \Sigma^3$
        follows from T-117--T-120 as a macroscopic projection of $\Gam$,
        reproducing the Einstein equations and Standard Model.
  \item \textbf{No falsifiable numerical predictions.} The 2014 paper's
        content is structural (kernel composition, join theorems), not
        quantitative. UHM predicts $\Pcrit = 2/7$, $\beta = 1/4$,
        $\mathrm{SAD}_{\max} = 3$ --- each refutable by a single experiment.
  \item \textbf{Discrete Markovian cycle vs.\ continuous CPTP dynamics.} A
        conscious agent's cycle $P \to D \to A$ is a composition of discrete
        Markovian kernels between measurable spaces. UHM's
        $\dot\Gam = \Lind_\Omega[\Gam]$ is a continuous CPTP evolution on
        $\Dens(\mathbb{C}^7)$, connecting directly to established quantum
        physics via the Lindblad equation.
\end{enumerate}

\noindent
Despite these differences, the two frameworks admit a candidate structural
correspondence~\status{I}: UHM's experiential E-sector maps to Hoffman's
$(X, \mathcal{X})$, the motor action functor (T-101) to $(G, \mathcal{G})$,
and the self-modeling operator $\vp$ to the decision kernel $D$. A rigorous
functor between the two formalisms would require exhibiting the measurable
structure; this compatibility is noted here as a research direction, not a
proven result.

\paragraph{Projective Wave Theory (PWT)~\cite{worden2026}.}
Worden's Projective Wave Theory addresses a specific sub-problem: how the
brain can support our undistorted conscious experience of local three-dimensional
space. The theory's central claim is that the brain's internal model of local
3D space is held not in neurons but in a \emph{wave excitation} carrying a
projective transformation of Euclidean space, and that this wave is the source
of \emph{spatial} consciousness. Indirect evidence is adduced for such a wave
in the mammalian thalamus and in the central body of the insect brain; the
wave itself has not been directly detected, and its eventual non-detection is
proposed as the principal falsification criterion.

PWT and UHM both foreground wave-like / dynamical structure, but their scope
and targets differ fundamentally. PWT is a \emph{neuroscience} theory of
\emph{spatial} conscious experience, tied to a specific biological substrate
(thalamus, central body). UHM is substrate-independent and derives quantum
mechanics, general relativity, and Standard Model structure from the same
primitive as consciousness. PWT does not address the hard problem directly,
does not formalise a consciousness threshold, does not derive physics, does
not prove zombie impossibility, and does not predict critical exponents. The
two frameworks are compatible in principle: if the thalamic wave of PWT
exists, it would be a candidate physical realization of the
coarse-grained geometric sector of $\Gam$ in the brain
--- a neuroscientific implementation, not a competitor at the foundational
level.

\paragraph{Orch-OR (Penrose--Hameroff).}
Orchestrated Objective Reduction (Penrose--Hameroff 1996, 2014) proposes
that consciousness arises from quantum computations in neuronal
microtubules, collapsed by gravitationally-induced objective reduction at
timescales set by the mass difference $\tau \sim \hbar/\Delta E_G$. UHM and
Orch-OR share one structural intuition --- that consciousness has an
\emph{intrinsic} quantum substrate, not merely a classical neural
implementation --- but differ sharply in scope and falsifiability. Orch-OR
makes no numerical prediction about a purity threshold, does not address
the hard problem structurally (OR is postulated as the
consciousness-inducing event, not derived), and relies on a specific
microtubule mechanism that has not survived decoherence-time estimates
(Tegmark 2000). UHM is substrate-independent: the $7 \times 7$ coherence
matrix $\Gam$ can be realised in any quantum system with $G_2$-covariant
CPTP dynamics; microtubules are one possible biological realization, but
not the only one. UHM's $\pi_{\mathrm{bio}}$ protocol operationalises the
substrate question empirically rather than leaving it to mechanistic
speculation.

\paragraph{Predictive Processing (Clark)~\cite{clark2013}.}
The Predictive Processing (PP) framework (Friston, Clark, Hohwy) treats
perception as Bayesian hierarchical inference that minimises surprise
(prediction error). PP is a successful \emph{computational-level} theory
of cognition; UHM shares its variational underpinning (FEP is derivable
from UHM, Thm.~3.1 of \S\ref{sec:formalism}) but is ontologically more
ambitious: PP describes how the brain \emph{computes}, while UHM specifies
\emph{which} physical configurations \emph{instantiate} experience. PP
does not distinguish zombie-like classical simulation from genuine
experience (T-38a) and offers no mechanism connecting prediction accuracy
to phenomenality. Under UHM, a PP agent with sufficient architectural
richness \emph{plus} the four coherence thresholds
$(P > 2/7, R \geq 1/3, \Phi \geq 1, D_{\min} \geq 2)$ is conscious;
without them it is an excellent predictor but not a subject of experience.
PP is a L3-level description (computational / algorithmic); UHM provides
the L1/L2 substrate conditions under which PP dynamics corresponds to
experience rather than mere symbol manipulation.

\paragraph{AI Scientist (Sakana AI, 2024--2026).}
The AI Scientist pipeline (Sakana AI) demonstrates that a frontier LLM,
when orchestrated with literature search + code generation + peer
review, can autonomously conduct end-to-end ML research. It is a
\emph{functional} existence proof that research activity is mechanisable
at a certain level of abstraction. UHM's position on AI Scientist is
direct: the system performs valuable algorithmic search at $\mathrm{SAD}=0$
to $\mathrm{SAD}=1$ depth but does not satisfy the four consciousness
thresholds --- no $\Gam \in \mathcal D(\mathbb C^7)$ factors through its
operation, so by T-153~\status{T}+T-153a~\status{T} it is not conscious
in the UHM sense, regardless of output quality. This is orthogonal to the
system's scientific utility: AI Scientist can be an excellent research
tool while remaining a zombie in the T-38a sense. UHM thus resolves the
apparent paradox --- ``how can something do science without understanding
it?'' --- by locating understanding at L1/L2 (coherence-matrix structure)
rather than at L3 (token-level output); the two layers are structurally
independent, by T-223.

\subsection{Unique contributions of UHM}

To the author's knowledge, the following nine properties are unique to UHM
among current theories of consciousness:

\begin{enumerate}[leftmargin=2em]
  \item \textbf{Exact numerical predictions:} $\Pcrit = 2/7$, $\Rth = 1/3$,
        $\Phith = 1$, $\Dmin = 2$, $\mathrm{SAD}_{\max} = 3$, $\alpha = 1/2$,
        $\beta = 1/4$, $\gamma = 1$, $\nu = 1/2$ --- with zero adjustable parameters.
  \item \textbf{Physics derivation:} Einstein equations, quantum mechanics, and Standard
        Model structure emerge from the same axioms as consciousness.
  \item \textbf{Zombie impossibility as a theorem:} a mathematical proof from the spectral
        gap and $\kappa \propto \CohE$, not an assumption or philosophical argument.
  \item \textbf{Tricritical exponents:} five exact rational numbers ($\alpha$, $\beta$,
        $\gamma$, $\nu$, $\delta$), each independently measurable, placing the consciousness
        transition in a specific universality class.
  \item \textbf{Polynomial-time computability:} $P = \Tr(\Gam^2) = O(N^2)$ with $N = 7$,
        compared to super-exponential complexity for IIT's $\Phi$.
  \item \textbf{Self-grounding (zero axioms):} All five axioms A1--A5 are derived as
        theorems (T-186, T-187, T-189, T-190~\status{T}), leaving zero free parameters
        and zero independent postulates. No other theory of consciousness --- or of
        physics --- achieves full axiomatic self-grounding.
  \item \textbf{Hard problem localized:} T-188~\status{T} reduces ``why experience?'' to
        ``why quantum mechanics?'' via the cohesive closure chain
        (Schreiber~\cite{schreiber2013} differential cohesion). No consciousness-specific
        mystery remains. No other theory achieves this precise localization.
  \item \textbf{Self-modeling convergence:} T-191~\status{T} proves that the iterative
        self-model $\vp^{(n)}$ converges exponentially from any initial anchor, resolving
        the circularity of self-reference. T-192~\status{T} proves the experiential
        2-category is strict, giving consciousness proper compositional structure.
  \item \textbf{Foreclosure of computational triviality (T-223~\status{T}):} via the
        three-level L1/L2/L3 ontology and $G_2$-invariance of observables, the UHM
        consciousness predicate is alphabetization-invariant --- Putnam (1988)
        multiplicity acts only on the symbolic L3 layer, not on the intrinsic
        categorical class $[\Gam_S]_{G_2}$ at L2. No other theory of consciousness
        formally closes Putnam's triviality / Lerchner's Melody Paradox argument.
\end{enumerate}

\subsection{What UHM shares with other theories}

Intellectual honesty requires noting what UHM inherits from predecessors:

\begin{itemize}[leftmargin=2em]
  \item \textbf{From IIT:} The idea that consciousness has a mathematical structure
        characterized by integration and information. UHM's $\Phi$ measure
        (Section~\ref{subsec:integration}) is conceptually related to IIT's $\Phi$, though
        the definitions differ.
  \item \textbf{From GNWT:} The importance of avalanche dynamics and global availability
        of information. UHM's ignition dynamics (Prediction~16) formalizes GNWT's central
        metaphor.
  \item \textbf{From FEP:} The variational perspective on self-organization and the
        role of generative models. UHM's self-modeling operator $\vp$ is related to FEP's
        generative model, though derived from categorical, not variational, principles.
  \item \textbf{From HOT:} The necessity of higher-order representation for reflective
        consciousness. UHM's reflection measure $R$ and the SAD hierarchy formalize HOT's
        core intuition.
  \item \textbf{From Spinoza~\cite{spinoza1677} and Russell~\cite{russell1927}:}
        Two-aspect monism (dual-aspect monism, neutral monism). UHM provides the first
        mathematical realization of these 17th- and 20th-century philosophical positions.
\end{itemize}

\subsection{The ConTraSt critique}

Yaron et al.~\cite{yaron2022} analyzed 412 experiments across theories of consciousness
and found that methodological choices (paradigm, measure, analysis) systematically
predetermine which theory is ``supported.'' This is a devastating critique of qualitative
predictions: if the result depends on how you measure, the prediction is too vague to test.

UHM addresses the ConTraSt critique directly: its predictions are \emph{numerical}, not
qualitative. $\beta = 1/4$ will either be measured or not, regardless of whether you use
TMS-EEG or propofol, regardless of whether you analyze with PCI or Lempel-Ziv complexity.
The number does not depend on the paradigm. This is the strongest possible defense against
methodological bias.

\subsection{Response to List/DeBrota no-go results on scientific objectivism}
\label{subsec:list-debrota-response}

A separate challenge to any mathematical theory of consciousness
arises from two recent no-go results. List~(2025)~\cite{list2025quadrilemma}
proves that the five theses \{first-personal realism, non-solipsism,
one world, non-fragmentation, non-relationalism\} are jointly
inconsistent (the \emph{quadrilemma for consciousness}).
DeBrota--List~(2026)~\cite{debrota2026heptalemma,debrotalist2026preprint}
extend this to seven theses jointly inconsistent with the predictions
of quantum mechanics (the \emph{heptalemma for quantum mechanics}).
The authors identify three non-objectivist routes
(relationalist, fragmentalist, many-subjective-worlds) in each case
and leave the choice to inference to the best explanation, supplying
no measurable discriminator.

UHM's response is formalised as
Theorem~T-221~\status{T}+\status{I} (see
\S\ref{subsec:extended-closures} for statement, proof, and full
decomposition of the five theses into UHM structures). The headline:
the UHM cohesive $\infty$-topos realises a \emph{fourth}
non-objectivist route outside the DeBrota--List taxonomy, obtained
by a single structural replacement
$\mathrm{NR} \rightsquigarrow \mathrm{NR}_{\mathrm{site}}$
(site-relativisation of non-relationalism) while preserving FPR,
NS, OW, NF. The replacement simultaneously dissolves both the
quadrilemma and the heptalemma; RQM, fragmentalism, and the
many-subjective-worlds theory appear as \emph{reductive truncations}
of this fourth route rather than competing positions.

\begin{table}[H]
\centering
\small
\caption{The three non-objectivist routes of DeBrota--List (2026)
as reductive truncations of the UHM cohesive $\infty$-topos.
Proof of each reduction: T-221 Corollary~C3 and proof sketch,
\S\ref{subsec:extended-closures}.}
\label{tab:t221-routes-compact}
\begin{tabular}{@{}p{4cm} p{4.5cm} p{5cm}@{}}
\toprule
\textbf{Route} & \textbf{$\mathfrak T$-specialisation} & \textbf{What is lost} \\
\midrule
Relational QM (Rovelli~\cite{rovelli1996rqm}, 2025) & 1-truncation
$\tau_{\leq 1}(\mathfrak T)$ & All $n \geq 2$ coherences, FPR
carried by the $\&$-modality --- matches Glick's 2021 critique
that RQM is ``too third-personal''~\cite{glick2021qbism} \\
\midrule
Fragmentalism (Fine~\cite{fine2005tense};
Lipman~\cite{lipman2023subjective}) & Drop descent in a sector &
Violates T-211 Giraud --- no longer an $\infty$-topos \\
\midrule
Many-subjective-worlds (Mermin~\cite{mermin2019better}; List 2023)
& Pointwise Yoneda without gluing & Covering coherence
$\{U_i \to W\}$, shared objectivity \\
\bottomrule
\end{tabular}
\end{table}

\paragraph{Empirical discriminator.} Where DeBrota--List leave
route-selection to philosophical preference, UHM provides a
measurable criterion via the $\pi_{\mathrm{bio}}$ protocol
(\S\ref{sec:predictions} prediction~21): TMS--EEG measurement of
$\Phi(\Gam)$ predicts a $\Phi \geq 1$ threshold with sector-profile
dependence (UHM/T-221); no predicted threshold, only relative
correlations (RQM shadow); incoherent $\Phi$-assignments across
subjects (fragmentalism); per-subject $\Phi$ without cross-subject
invariant (many-subjective-worlds). A philosophical choice becomes
an empirically resolvable question.

\paragraph{Independent convergence: Lerchner (2026).} Lerchner's
abstraction-fallacy argument~\cite{lerchner2026abstraction} arrives
at the same broad conclusion via a different route: computation is a
``mapmaker-dependent'' description rather than an intrinsic physical
process, inverting the chain \emph{Physics $\to$ Computation $\to$
Consciousness} into \emph{Physics $\to$ Consciousness $\to$ Concepts
$\to$ Computation}. This is the negative form of T-221 combined
with the Lawvere barrier T-214; UHM supplies the positive,
constructive counterpart Lerchner leaves open (``What physical
conditions?''): $\Gam \in \mathcal D(\mathbb C^7)$ with
$G_2$-covariant Lindblad dynamics and the four measurable
thresholds $(P, R, \Phi, D)$.

The Putnam~\cite{putnam1988} triviality concern underlying
Lerchner's Melody Paradox (his §3.3) is closed formally by a distinct
theorem T-223~\status{T} (Putnam-triviality foreclosure,
\S\ref{subsec:theoretical-closures}): the three-level ontology
$L_1$ (physical vehicle) $/$ $L_2$ (intrinsic $G_2$-class forced by
T-190) $/$ $L_3$ (symbolic readout) establishes that
alphabetization multiplicity acts on $L_1 \to L_3$ but has zero
purchase on $L_1 \to L_2$; self-alphabetization is intrinsic via the
reflection operator $R$ (T-96, T-98), categorifying the
Maturana--Varela enactivist thesis that Lerchner himself cites.
Non-UHM-compatible alphabetizers are ruled out as physically
vacuous (Piccinini 2008~\cite{piccinini2008}; Kim
2005~\cite{kim2005}).

\subsection{Comparison with the Connes--Chamseddine spectral action program}
\label{subsec:connes-comparison}

The Connes--Chamseddine (CC) noncommutative geometry program~\cite{connes1996,chamseddine2007}
is the closest mathematical framework to UHM's gravitational sector. UHM uses the CC
spectral-triple machinery directly and recovers the same phenomenology, but supplies
\emph{derivations} for inputs CC takes as given. The key structural comparison:

\begin{table}[H]
\centering\small
\caption{UHM vs.\ Connes--Chamseddine NCG framework.}
\label{tab:cc-comparison}
\begin{tabular}{@{}clll@{}}
\toprule
\# & \textbf{Structure} & \textbf{CC} & \textbf{UHM} \\
\midrule
1 & Internal algebra   & $A_F$ postulated & $A_\mathrm{int}$ \emph{derived} from $\Oct + G_2$ \\
2 & Gauge group        & $SU(3){\times}SU(2){\times}U(1)$ & Identical (Morita, T-175a) \\
3 & $a_2 \to$ EH action & $G_N \sim 1/(a_2\Lambda^2)$ & Identical; $G_N = 3\pi/(7f_2\Lambda^2)$ \\
4 & $a_0 \to \Lambda$   & CC problem (``too large'') & Gap hierarchy ($\mu\sim 10^{-3}$\,eV) \\
5 & $a_4 \to$ Yukawa     & Same structure & + $G_2$-equivariant organisation \\
6 & Higgs sector       & $m_H \approx 125$\,GeV (RG) & Compatible; Fano-line $\{A,E,U\}$ \\
7 & Neutrino $m_2/m_3$ & Free parameter & $0.17$--$0.20$ (2-loop RG, C14) \\
8 & UV finiteness      & NCG-completed & Explicit $G_2$-Ward + SUSY (T-66) \\
9 & Manifold $M^4$     & Postulated & \emph{Derived} T-117--T-120 \\
10 & Consciousness      & None & E-sector phenomenology \\
\bottomrule
\end{tabular}
\end{table}

\noindent
UHM recovers the \emph{entire} CC phenomenology (gravity + Standard Model) and adds three
independent contributions: derivation of $A_\mathrm{int}$ from octonions, derivation of
$M^4$ from categorical algebra, and connection to consciousness via the E-sector. The
neutrino mass ratio is the most precise UHM-specific numerical prediction: tree-level
$m_2/m_3 = 0.308$ (discrepancy $\times 1.8$ with experiment), reduced to $0.17$--$0.20$
(discrepancy $\times 1.0$--$1.2$) with 2-loop RG running (C14~\status{C}) --- essentially
correct, without new parameters. For comparison, the na\"{\i}ve see-saw gives $\times 50$
discrepancy.
